% 第12章 ブラック・ショールズ・モデル(資料12 + BS公式導出補足)
\chapter{ブラック・ショールズ・モデル}

最終章では，連続時間モデルによるデリバティブの価格付け，すなわち
\textbf{ブラック・ショールズ(BS)モデル}を学ぶ。離散時間の二項モデル(CRR)から
連続時間モデルへ移行し，ヨーロピアン・オプションの解析的な価格公式
(Black-Scholesの公式)を導出する。また，デリバティブのリスク管理に使われる
\textbf{グリークス} (Greeks) と，通貨・金利デリバティブへの応用の考え方を扱う。

\section{ブラック・ショールズ(BS)モデル}

\subsection{前提}

\begin{itemize}
\item 完全市場(第9章参照)を仮定する。
\item 安全資産の利回り $r$($=$無リスク金利)は期間を通して一定とする。
\end{itemize}
さらに，簡単化のために以下を仮定する(これらの仮定は緩和可能である)。
\begin{itemize}
\item 市場には，1つのリスク性資産($=$原資産)と無リスク資産だけが存在する
(ということは，複製はこれら2資産を使って考える)。
\item リスク性資産に配当はない。
\end{itemize}
以下では，\textbf{リスク中立確率のもとで}，連続時間モデルで考える(連続時間
なので，ほんの一瞬の間の変化を考える)。

\subsection{モデルの定式化}

Black-Scholes(-Merton) モデルでは，微小時間 $\Delta t$ における瞬間的な収益率が，
\[
\text{期待値 } r\Delta t \quad
(\text{リスク中立確率下の議論なので } r \neq \mu),\qquad
\text{分散 } \sigma^2\Delta t
\]
の正規分布に従う，と仮定する。$\Delta$ で増分を表すと，
\[
\frac{\Delta S(t)}{S(t)} \sim N(r\Delta t,\ \sigma^2\Delta t),
\qquad \Delta S(t) = S(t+\Delta t) - S(t).
\]
無限小時間 $dt$ における式として表現すると，
\[
\frac{\dd S(t)}{S(t)} \sim N(r\,\dd t,\ \sigma^2\dd t),
\qquad \dd S(t) = S(t+\dd t) - S(t).
\]
しかも，瞬間的な収益率は時系列的に独立と仮定する。

これを，以下のように表現する：
\begin{equation}
\frac{\dd S(t)}{S(t)} = r\,\dd t + \sigma\,\dd W(t).
\label{eq:bs-sde}
\end{equation}
\begin{itemize}
\item 左辺は，無限小時間 $dt$ における瞬間的な収益率を表す。
\item 右辺第1項は，期待収益率 $r$ による $dt$ 間の増分を表す。
\item 右辺第2項は，$dt$ 間の収益率の確率的な増分(変動)を表す。変動の大きさは
$\sigma$ で調整する。
\end{itemize}

\begin{definition}[標準ブラウン運動]
$W(t)$ は\textbf{標準ブラウン運動}と呼ばれている確率過程で，次の性質をもつ。
\begin{itemize}
\item $\dd W(t) = W(t+\dd t) - W(t)$：標準ブラウン運動の無限小増分。
\item 無限小増分 $\dd W(t)$ は正規分布 $N(0,\,\dd t)$ に従う
($\Rightarrow \E[\dd W(t)] = 0$)。
\item 互いに背反な時間区間の無限小増分は，独立。
\item $W(0) = 0$， a.s.
\end{itemize}
\end{definition}

\begin{example}[問1]
式\eqref{eq:bs-sde}が $\dd S(t)/S(t) \sim N(r\,\dd t, \sigma^2\dd t)$ を満たす
ことを確認せよ。

\textbf{解。}(1) 期待値と分散を計算する：
\[
\E\Bigl[\frac{\dd S(t)}{S(t)}\Bigr] = r\,\dd t + \sigma\E[\dd W(t)] = r\,\dd t,
\qquad
\Var\Bigl[\frac{\dd S(t)}{S(t)}\Bigr] = \sigma^2\Var[\dd W(t)] = \sigma^2\dd t.
\]
(2) 右辺 $r\,\dd t + \sigma\,\dd W(t)$ は，正規分布に従う確率変数
$\dd W(t)$ の1次式なので正規分布に従う。よって
$\dd S(t)/S(t) \sim N(r\,\dd t, \sigma^2\dd t)$ である。
\end{example}

\begin{example}[問2]
では，$W(t)$ はどのような分布に従うか?

\textbf{解。}(1)
\[
W(t) = \int_0^t \dd W(s)
= \lim_{n\to\infty}\sum_{i=1}^{n}
\Bigl[W\Bigl(\frac{it}{n}\Bigr) - W\Bigl(\frac{(i-1)t}{n}\Bigr)\Bigr]
\]
と表現して，互いに背反な時間区間の増分の和に分解する。各増分は独立に
$N(0, t/n)$ に従う。(2) 正規分布の再生性(独立な正規変数の和は正規分布)を
使うと，$W(t) \sim N(0,\,t)$ が得られる。
\end{example}

\vspace{18pt}

\subsection{単純な積分ではうまくいかない}

式\eqref{eq:bs-sde}はリスク性資産価格の\emph{瞬間的な増分}の分布を示す式で
あって，将来時点 $t$ の価格 $S(t)$ の分布を示してはいない。単純に積分しても
有益な情報は出てこない：
\[
\dd S(t) = rS(t)\,\dd t + \sigma S(t)\,\dd W(t)
\]
\[
\int_0^T \dd S(t) = r\int_0^T S(t)\,\dd t + \sigma\int_0^T S(t)\,\dd W(t)
\]
\[
S(T) = S(0) + r\int_0^T S(t)\,\dd t + \sigma\int_0^T S(t)\,\dd W(t).
\]
被積分関数に確率過程 $S(t)$ が含まれたままなので，$S(T)$ の分布が上の式から
直接得られることはない。

そこで，算術収益率でなく，\textbf{対数収益率}で考えてみる(結果から言うと，
この方針をとるとうまくいく)。瞬間的な対数収益率は
$\dd\log S(t) = \log S(t+\dd t) - \log S(t)$ である。$\dd\log S(t)$ はどのような
分布に従うか。結果を先取りすると，$\dd\log S(t) \sim N(r\,\dd t, \sigma^2\dd t)$
には\emph{ならない}。やりたいことは，$f(t,S) = \log S$ の微小増分
$\dd f(t,S(t))$ を求めることである。このように，微小増分 $\dd S(t)$ の分布が
既知のとき，微小増分 $\dd f(t, S(t))$ の分布がどうなるかを考える時のツールが
\textbf{伊藤の公式}である。詳しく言うと，ある確率過程が与えられたとき，その
確率過程を含む汎関数(関数の関数)がどのように書けるか，を示す公式である。

\subsection{伊藤の公式}

$\dd S(t) = S(t)(r\,\dd t + \sigma\,\dd W(t))$ は既知とする。二変数関数
$f(t, S(t))$ を考え，無限小時間 $dt$ における微小増分 $\dd f(t,S(t))$ を考える。
2変数関数のテーラー展開を使うと，
\begin{equation}
\dd f(t,S(t)) = \frac{\partial f}{\partial t}\dd t
+ \frac{\partial f}{\partial S}\dd S(t)
+ \frac{1}{2}\frac{\partial^2 f}{\partial t^2}(\dd t)^2
+ \frac{\partial^2 f}{\partial t\,\partial S}\dd t\,\dd S(t)
+ \frac{1}{2}\frac{\partial^2 f}{\partial S^2}(\dd S(t))^2 + \cdots.
\end{equation}
ここで $\dd S(t) = S(t)(r\,\dd t + \sigma\,\dd W(t))$ を代入して整理すると，
\begin{align}
\dd f(t,S(t))
&= \frac{\partial f}{\partial t}\dd t
+ \frac{\partial f}{\partial S}S(t)\bigl(r\,\dd t + \sigma\,\dd W(t)\bigr)
+ \frac{1}{2}\frac{\partial^2 f}{\partial S^2}S(t)^2
\bigl(r\,\dd t + \sigma\,\dd W(t)\bigr)^2 + \cdots\notag\\
&= \Bigl(\frac{\partial f}{\partial t} + rS(t)\frac{\partial f}{\partial S}\Bigr)
\dd t + \sigma S(t)\frac{\partial f}{\partial S}\dd W(t)
+ \frac{1}{2}\sigma^2 S(t)^2\frac{\partial^2 f}{\partial S^2}(\dd W(t))^2
+ (\text{高次項}).
\end{align}
1次より高いオーダー($(\dd t)^2$， $\dd t\,\dd W(t)$ 等)を無視すると，最後に
$(\dd W(t))^2$ の項が残る。

\begin{keypoint}[$(\dd W(t))^2 = \dd t$：このページが肝心!]
重要：$\dd W(t)$ のオーダーは $(\dd t)^{0.5}$ である。実際，
\[
\Var[\dd W(t)] = \E\bigl[(\dd W(t) - \E[\dd W(t)])^2\bigr]
= \E[(\dd W(t))^2] = \dd t,
\]
であり，さらに($N(0,\dd t)$ の4次モーメント $3(\dd t)^2$ を使うと)
\begin{align*}
\Var[(\dd W(t))^2]
&= \E\bigl[((\dd W(t))^2 - \E[(\dd W(t))^2])^2\bigr]
= \E\bigl[((\dd W(t))^2 - \dd t)^2\bigr]\\
&= \E[(\dd W(t))^4] - 2\,\dd t\,\E[(\dd W(t))^2] + (\dd t)^2
= 3(\dd t)^2 - (\dd t)^2 = 2(\dd t)^2.
\end{align*}
よって $(\dd W(t))^2$ は平均 $\dd t$，分散 $2(\dd t)^2$ である。分散は $\dd t$ の
2次項なので無視すると，\textbf{$(\dd W(t))^2 = \dd t$ とみなせる}。
\end{keypoint}

上記を使い，$\dd t$ と $\dd W(t)$(一次)より高次の項を無視すると，
\begin{equation}
\dd f(t,S(t))
= \Bigl(\frac{\partial f}{\partial t} + rS(t)\frac{\partial f}{\partial S}
+ \frac{1}{2}\sigma^2S(t)^2\frac{\partial^2 f}{\partial S^2}\Bigr)\dd t
+ \sigma S(t)\frac{\partial f}{\partial S}\,\dd W(t).
\label{eq:ito}
\end{equation}
この式を\textbf{伊藤の公式}という。

\subsection{対数価格の確率微分方程式と $S(T)$ の分布}

$f(t,S(t)) = \log S(t)$ として伊藤の公式\eqref{eq:ito}を用いると，
$\partial f/\partial t = 0$， $\partial f/\partial S = 1/S$，
$\partial^2 f/\partial S^2 = -1/S^2$ より
\begin{equation}
\dd\log S(t)
= \Bigl(r - \frac{1}{2}\sigma^2\Bigr)\dd t + \sigma\,\dd W(t)
\label{eq:dlogS}
\end{equation}
となる。すなわち，
$\dd S(t)/S(t) = r\,\dd t + \sigma\,\dd W(t)$
($\dd S(t) = rS(t)\dd t + \sigma S(t)\dd W(t)$)と違い，
\eqref{eq:dlogS}は\textbf{右辺の $\dd t$ 項の係数と $\dd W(t)$ 項の係数が定数
(確定関数)}である。そのため，両辺を時間積分した結果が明確に書ける。

そこで，以後は Black-Scholesモデルとして
\begin{equation}
\dd\log S(t) = \Bigl(r - \frac{1}{2}\sigma^2\Bigr)\dd t + \sigma\,\dd W(t)
\end{equation}
を考える。両辺積分すると
\[
\log\frac{S(T)}{S(0)} = \Bigl(r - \frac{1}{2}\sigma^2\Bigr)T + \sigma W(T)
\]
となるので，
\begin{equation}
S(T) = S(0)\exp\Bigl\{\Bigl(r - \frac{1}{2}\sigma^2\Bigr)T + \sigma W(T)\Bigr\},
\qquad T \geq 0.
\label{eq:ST}
\end{equation}

\begin{keypoint}[BSモデルの特徴]
\begin{itemize}
\item $S(0) > 0$ ならば，常に $S(T) > 0$。
\item $\log S(T)$ は正規分布
$N\bigl(\log S(0) + (r-\sigma^2/2)T,\ \sigma^2 T\bigr)$ に従う。
$\Rightarrow$ $S(T)$ は\textbf{対数正規分布}に従う。
\end{itemize}
\end{keypoint}

\section{BSモデルによるオプション価格公式の導出}

\subsection{方針}

ヨーロピアン・コール・オプションとプット・オプションの価格式が，BSモデルでは
どのように表現されるかを導出したい。

(過去の資料の方針)手掛かりは，多期間二項モデルで得られたCRRの公式と，前述の
BSモデルに関する数式である。離散時間モデル(CRR)から連続時間モデル(BS)へ
移行する。そこで使うロジックは：満期を固定し，期間分割数を $n\to\infty$ とした
とき，多期間二項モデルがBSモデルに収束するようにパラメータを設定すれば，CRRの
公式はBSモデルの下でのオプション価格式に収束するはず，というものである。

(この資料の方針)せっかく\textbf{リスク中立化法}を学んだので，それを使用する。
すなわち，「リスク中立確率の下における将来キャッシュフローの割引現在価値の
期待値が，証券(デリバティブ)の現在の価格になる」ことを使う。

\subsection{連続時間モデルのリスク中立化法}

これを連続時間モデルで表現すると，金利 $r$ が一定のとき，
\begin{equation}
C_0 = \E^Q\bigl[e^{-rT}f(S(T))\bigr] = e^{-rT}\,\E^Q\bigl[f(S(T))\bigr].
\label{eq:rn-valuation}
\end{equation}
ここで，$T$：満期，$S(t)$：時刻 $t$ の原資産価格，$f(\cdot)$：デリバティブの
ペイオフ関数である。例えばヨーロピアンコールオプションでは
$f(S(T)) = \max(S(T)-K,\,0)$ である。

\subsection{コールの価格式の導出}

原資産価格 $S(t)$，満期 $T$，行使価格 $K$ のヨーロピアン・コール・オプションの
現在価値を求める。リスク中立化法を適用すると，
\begin{equation}
C_0 = e^{-rT}\E^Q[f(S(T))] = e^{-rT}\E^Q[\max(S(T)-K,\,0)].
\label{eq:call-rn}
\end{equation}
以後の式展開は，配布資料「Black-Scholesモデルによるヨーロピアン・コール・
オプション価格式」に基づく(以下，本節に全文を収録する)。

\paragraph{準備：$\log S(T)$ の密度関数。}
Black-Scholesモデルでは満期 $T$ の価格 $S(T)$ は，式\eqref{eq:ST}より
\[
S(T) = S\exp\Bigl\{\Bigl(r-\frac{\sigma^2}{2}\Bigr)T + \sigma Z^Q(T)\Bigr\}
\]
である。ただし $S = S(0)$，$Z^Q(t)$ は確率 $Q$ の下における標準ブラウン運動，
$r, \sigma$ は定数である。これより
\[
X(T) = \log S(T) = \log S + \Bigl(r-\frac{\sigma^2}{2}\Bigr)T + \sigma Z^Q(T)
\]
なので，$X(T)$ は平均
\[
\mu = \log S + \Bigl(r - \frac{\sigma^2}{2}\Bigr)T,
\]
分散 $\sigma^2 T$ の正規分布に従い，$X(T)$ の密度関数は
\begin{equation}
f_X(x) = \frac{1}{\sqrt{2\pi T}\,\sigma}\,
e^{-\frac{(x-\mu)^2}{2\sigma^2T}}
\end{equation}
である。そこで\eqref{eq:call-rn}を
\begin{align}
C_0 &= e^{-rT}\E^Q\bigl[\max(e^{X(T)} - K,\,0)\bigr]
= e^{-rT}\int_{-\infty}^{\infty}\max(e^x - K,\,0)\,f_X(x)\,\dd x
\end{align}
とすることにする。

\paragraph{積分領域の限定。}
この積分は $\max(\cdot,\cdot)$ を含むため，積分に寄与するのは
\[
\{x : \max(e^x - K, 0) > 0\} = \{x : e^x > K\} = \{x : x > \log K\}
\]
の領域だけである。よって
\begin{align}
C_0 &= e^{-rT}\int_{\log K}^{\infty}(e^x - K)\,
\frac{1}{\sqrt{2\pi T}\sigma}e^{-\frac{(x-\mu)^2}{2\sigma^2T}}\dd x
= e^{-rT}I_1 - Ke^{-rT}I_2,
\label{eq:c0-i1i2}
\end{align}
\begin{equation}
I_1 = \int_{\log K}^{\infty}e^x\,
\frac{1}{\sqrt{2\pi T}\sigma}e^{-\frac{(x-\mu)^2}{2\sigma^2T}}\dd x,
\qquad
I_2 = \int_{\log K}^{\infty}
\frac{1}{\sqrt{2\pi T}\sigma}e^{-\frac{(x-\mu)^2}{2\sigma^2T}}\dd x
\end{equation}
と書き直す。

\paragraph{$I_2$ の計算。}
変数変換
\[
y = \frac{x - \mu}{\sigma\sqrt{T}}
\]
をすると，
\begin{equation}
I_2 = \int_{\frac{\log K - \mu}{\sigma\sqrt{T}}}^{\infty}
\frac{1}{\sqrt{2\pi}}e^{-\frac{y^2}{2}}\dd y
= 1 - \Phi\Bigl(\frac{\log K - \mu}{\sigma\sqrt{T}}\Bigr)
= \Phi\Bigl(\frac{\mu - \log K}{\sigma\sqrt{T}}\Bigr).
\end{equation}
ただし，$\Phi(\cdot)$ は標準正規分布の分布関数である。

\paragraph{$I_1$ の計算。}
まず指数部を平方完成する：
\begin{align}
I_1 &= \int_{\log K}^{\infty}\frac{1}{\sqrt{2\pi T}\sigma}
e^{-\frac{(x-\mu)^2 - 2\sigma^2Tx}{2\sigma^2T}}\dd x
= \int_{\log K}^{\infty}\frac{1}{\sqrt{2\pi T}\sigma}
e^{-\frac{x^2 - 2(\mu+\sigma^2T)x + \mu^2}{2\sigma^2T}}\dd x\notag\\
&= \int_{\log K}^{\infty}\frac{1}{\sqrt{2\pi T}\sigma}
e^{-\frac{(x-(\mu+\sigma^2T))^2 + \mu^2 - (\mu+\sigma^2T)^2}{2\sigma^2T}}\dd x
= e^{\frac{(2\mu+\sigma^2T)\sigma^2T}{2\sigma^2T}}
\int_{\log K}^{\infty}\frac{1}{\sqrt{2\pi T}\sigma}
e^{-\frac{(x-(\mu+\sigma^2T))^2}{2\sigma^2T}}\dd x.
\end{align}
変数変換
\[
u = \frac{x - (\mu + \sigma^2 T)}{\sigma\sqrt{T}}
\]
を行うと，
\begin{align}
I_1 &= e^{\frac{2\mu+\sigma^2T}{2}}
\int_{\frac{\log K - (\mu+\sigma^2T)}{\sigma\sqrt{T}}}^{\infty}
\frac{1}{\sqrt{2\pi}}e^{-\frac{u^2}{2}}\dd u\notag\\
&= e^{\frac{2\mu+\sigma^2T}{2}}
\Bigl(1 - \Phi\Bigl(\frac{\log K - (\mu+\sigma^2T)}{\sigma\sqrt{T}}\Bigr)\Bigr)\notag\\
&= e^{\frac{2\mu+\sigma^2T}{2}}
\Phi\Bigl(\frac{\mu + \sigma^2T - \log K}{\sigma\sqrt{T}}\Bigr)
\end{align}
となる。このままでもよいが，$\mu$ をもとの変数で表すと，
\begin{align}
\frac{\mu - \log K}{\sigma\sqrt{T}}
&= \frac{\log\frac{S}{K} + \bigl(r - \frac{\sigma^2}{2}\bigr)T}{\sigma\sqrt{T}},\\
\frac{\mu + \sigma^2T - \log K}{\sigma\sqrt{T}}
&= \frac{\log\frac{S}{K} + \bigl(r + \frac{\sigma^2}{2}\bigr)T}{\sigma\sqrt{T}},\\
\frac{2\mu + \sigma^2 T}{2}
&= \frac{2\log S + 2\bigl(r-\frac{\sigma^2}{2}\bigr)T + \sigma^2T}{2}
= \log S + rT
\end{align}
となる。

\paragraph{Black-Scholesのコール価格式。}
以上を\eqref{eq:c0-i1i2}に代入すると，ヨーロピアン・コール・オプションの価格式
\begin{align}
C_0 &= S\,\Phi\Biggl(\frac{\log\frac{S}{K}
+ \bigl(r+\frac{\sigma^2}{2}\bigr)T}{\sigma\sqrt{T}}\Biggr)
- Ke^{-rT}\,\Phi\Biggl(\frac{\log\frac{S}{K}
+ \bigl(r-\frac{\sigma^2}{2}\bigr)T}{\sigma\sqrt{T}}\Biggr)\notag\\
&= S\,\Phi\Biggl(\frac{\log\frac{S}{K}
+ \bigl(r+\frac{\sigma^2}{2}\bigr)T}{\sigma\sqrt{T}}\Biggr)
- Ke^{-rT}\,\Phi\Biggl(\frac{\log\frac{S}{K}
+ \bigl(r+\frac{\sigma^2}{2}\bigr)T}{\sigma\sqrt{T}} - \sigma\sqrt{T}\Biggr)
\label{eq:bs-call}
\end{align}
が得られる。

\begin{remark}
式\eqref{eq:bs-call}の第2項の $\Phi(\cdot)$ は，リスク中立確率 $Q$ の下で満期
$T$ においてコールが行使される確率である。つまり，第2項は「行使時に支払う
行使価格の期待割引価値(行使確率を織り込んだもの)」を意味している。
\end{remark}

\vspace{18pt}

\subsection{Black-Scholesの公式}

\begin{keypoint}[Black-Scholesの公式]
\begin{align}
V_0^C &= S\,\Phi(a_1) - Ke^{-rT}\,\Phi\bigl(a_1 - \sigma\sqrt{T}\bigr),\notag\\
V_0^P &= Ke^{-rT}\,\Phi\bigl(-a_1+\sigma\sqrt{T}\bigr) - S\,\Phi(-a_1),\notag\\
a_1 &= \frac{\log(S/K) + (r + \sigma^2/2)\,T}{\sigma\sqrt{T}}.
\label{eq:bs-formula}
\end{align}
ここまでで求めた，連続時間モデルにおけるヨーロピアン・コール・
オプションおよびヨーロピアン・プット・オプションの価格式を，
\textbf{Black-Scholesの公式}(Black-Scholesのオプション価格式)という。
\end{keypoint}

\paragraph{プット価格式の導出。}
プット・コール・パリティ($r$ 一定なので $v(t,T) = e^{-r(T-t)}$ となることに
注意)より，プットオプションの価格は
\begin{align}
V_0^P &= V_0^C - S + Ke^{-rT}
= S\,\Phi(a_1) - Ke^{-rT}\Phi(a_1-\sigma\sqrt{T}) - S + Ke^{-rT}\notag\\
&= S\bigl(\Phi(a_1) - 1\bigr)
+ Ke^{-rT}\bigl(1 - \Phi(a_1-\sigma\sqrt{T})\bigr)\notag\\
&= Ke^{-rT}\,\Phi\bigl(-a_1 + \sigma\sqrt{T}\bigr) - S\,\Phi(-a_1)
\label{eq:bs-put}
\end{align}
となる。コール・オプションと同様の計算(省略)をしても導出できる。

\subsection{BS公式による複製ポートフォリオ}

\begin{itemize}
\item ヨーロピアン・コール・オプションの複製ポートフォリオ：「複製ポート
フォリオはリスク性資産 $\Delta$ 単位，(一単位あたり1円の)無リスク資産
$B$ 単位からなる」とすると，価格式\eqref{eq:bs-call}から
\begin{equation}
\Delta = \Phi(a_1),\qquad
B = -Ke^{-rT}\,\Phi\bigl(a_1 - \sigma\sqrt{T}\bigr).
\end{equation}
\item ヨーロピアン・プット・オプションの複製ポートフォリオ：同様に，価格式
\eqref{eq:bs-put}から
\begin{equation}
\Delta = -\Phi(-a_1),\qquad
B = Ke^{-rT}\,\Phi\bigl(-a_1 + \sigma\sqrt{T}\bigr).
\end{equation}
\end{itemize}

\subsection{本源的価値と時間価値}

コール価格のグラフ($K=100$， $r=0.1$， $\sigma=0.2$， $T=1$ の例)を原資産価格の
関数として描くと，満期のペイオフ $\max(S-K,0)$(これを\textbf{本源的価値}と
いう)より上に位置する滑らかな曲線になる。オプション価格と本源的価値の差を
\textbf{時間価値}という。

\begin{itemize}
\item ヨーロピアン\emph{コール}オプションでは，時間価値は正になることが多い。
\item しかし，ヨーロピアン\emph{プット}オプションでは，ITMが深い(原資産価格が
十分低い)領域で，オプション価格が本源的価値を\emph{下回る}(時間価値が負に
なる)ことがある。
\end{itemize}

\subsection{BS公式の振る舞い}

BS公式は多くの変数に依存するので，振る舞いは複雑である。ここでは
$V_0^C = V_0^C(S, K, T, \sigma, r)$ について，一つ一つの変数の動きと価格の
動きをみる(証明・確認は各自で考えること)。

\begin{itemize}
\item 原資産価格 $S$：
$S \to 0$ で $V_0^C \to 0$，\quad $S\to\infty$ で $V_0^C \to\infty$。
\item 行使価格 $K$：
$K \to 0$ で $V_0^C \to S$，\quad $K\to\infty$ で $V_0^C \to 0$。
\item 満期 $T$：
$T\to 0$ で $V_0^C \to \max(S-K, 0)$
($S\geq K$ なら $S-K$ へ，$S<K$ なら0へ)，\quad
$T\to\infty$ で $V_0^C \to S$。
\item ボラティリティ $\sigma$：
$\sigma\to 0$ で $V_0^C \to \max(S - Ke^{-rT},\,0)$
($S\geq Ke^{-rT}$ なら $S - Ke^{-rT}$ へ，$S<Ke^{-rT}$ なら0へ)，\quad
$\sigma\to\infty$ で $V_0^C \to S$。
\item 無リスク金利 $r$：
$r\to 0$ で $V_0^C \to V_0^C(S,K,T,\sigma,0)$，\quad
$r\to\infty$ で $V_0^C \to S$。
\end{itemize}

\begin{remark}[参考：デリバティブ価格の求め方]
デリバティブ価格の求め方には，
\begin{enumerate}
\item リスク中立化法を使い，(将来キャッシュフローの割引現在価値の)期待値を
解析的に計算する。
\item リスク中立化法を使い，期待値をモンテカルロ・シミュレーションで計算する。
\item 二項モデルのような格子モデルで，期待値を計算する。
\item 価格関数が満たす偏微分方程式を解き，価格を求める。
\end{enumerate}
などがある。本講義では1と3だけを扱った。
\end{remark}

\vspace{18pt}

\section{デリバティブのリスク管理：グリークス (Greeks)}

\subsection{リスクヘッジとリスク感応度}

\begin{itemize}
\item \textbf{リスク・ヘッジ}とは：ここでいうリスクとは，将来発生しうる価格の
変動性のこと。そのようなリスクを回避・低減することをリスク・ヘッジという。
\item \textbf{リスク感応度}：リスクの源泉となる変数が1単位動いたときの価格
(or 損失額)の変化量。(偏)微分係数で与えることもある(グリークスはこの
偏微分係数のこと)。変数の価格への影響度を示す。
\end{itemize}

\subsection{グリークス}

オプション価格 $V$ をテーラー展開すると($t$ は時刻，$T$ は満期までの期間で，
$\dd T = -\dd t$)，
\begin{align}
&V(S+\dd S,\ K,\ T-\dd t,\ \sigma+\dd\sigma,\ r+\dd r)\notag\\
&\quad = V(S,K,T,\sigma,r)
+ \frac{\partial V}{\partial S}\dd S
- \frac{\partial V}{\partial T}\dd t
+ \frac{\partial V}{\partial\sigma}\dd\sigma
+ \frac{\partial V}{\partial r}\dd r
+ \frac{1}{2}\frac{\partial^2V}{\partial S^2}(\dd S)^2 + \cdots.
\end{align}
これより，オプション価格の変化は幾つかの偏微分を使って表現できる。
これらのギリシャ文字で表現されるリスク指標を\textbf{グリークス}という。
\begin{itemize}
\item まずは一次の項を考える。実務では，デルタ・ヘッジ(後出)は当たり前の
ように実施する。
\item 実務では，ベガ・リスク(後出)と若干の高次のリスクにも着目する。
\item 二次の項で考慮するのは主にガンマ(後出)である。
\end{itemize}
以下ではヨーロピアン・コール・オプション
$V_0^C = V_0^C(S,K,T,\sigma,r)$ を中心に具体的な式を出す。

\subsection{デルタ($\Delta$)，ガンマ($\Gamma$)}

\begin{definition}[デルタ]
\textbf{デルタ} (Delta) は，原資産価格 $S$ が微小量変化したときの，オプション
価格の変化率である。オプション価格 $V$ の原資産価格 $S$ に関する偏微分で表現
される。ヨーロピアン・コールでは
\begin{equation}
\Delta = \frac{\partial V_0^C}{\partial S} = \Phi(a_1).
\end{equation}
\end{definition}

これは複製ポートフォリオにおける原資産の保有量 $\Delta = \Phi(a_1)$ と同じで
ある(「デルタ」の別の意味)。実際には，オプション価格 $V$ は原資産価格 $S$ の
非線形関数であり，$\Delta$ は $S$ に依存して時々刻々変化する。ということは，
\emph{複製ポートフォリオも時々刻々変化する}。

\begin{definition}[ガンマ]
$\Delta$ の $S$ への依存関係を示すのが\textbf{ガンマ} (Gamma) である：
\begin{equation}
\Gamma = \frac{\partial\Delta}{\partial S}
= \frac{\partial^2 V_0^C}{\partial S^2}
= \frac{1}{S\sigma\sqrt{T}}\,
\left.\frac{\dd\Phi(x)}{\dd x}\right|_{x=a_1}
= \frac{\varphi(a_1)}{S\sigma\sqrt{T}}
\end{equation}
($\varphi$ は標準正規密度関数)。
\end{definition}

$K=100$， $r=0.1$， $\sigma=0.2$， $T=1$ の例では，$\Delta$ は原資産価格とともに
0から1へS字型に増加し，$\Gamma$ はATM付近で最大となる山型になる。

\begin{keypoint}[ガンマとリバランス]
ガンマはリバランスのタイミングの目安の一つになる。
\begin{itemize}
\item ガンマが小さければ，原資産価格が少し変化してもデルタはあまり変化しない
$\to$ リバランスはやや間隔を空けても大丈夫…かも?
\item ガンマが大きければ，原資産価格が少し変化するだけでデルタはかなり変化する
$\to$ リバランスは頻繁に行うほうがよい…かも?
\end{itemize}
これらはリバランスに要するコストとの兼ね合いで決める。
\end{keypoint}

\subsection{セータ($\Theta$)}

\begin{definition}[セータ]
\textbf{セータ}(シータ，Theta)は，満期 $T$ が微小量変化したときのオプション
価格の変化率(満期を固定したオプションの価格の，瞬間的な変化率)である：
\begin{align}
\Theta &= \frac{\partial V_0^C}{\partial T}
= \frac{\sigma S}{2\sqrt{T}}\,\left.\frac{\dd\Phi(x)}{\dd x}\right|_{x=a_1}
+ rKe^{-rT}\,\Phi\bigl(a_1 - \sigma\sqrt{T}\bigr)\notag\\
&= \frac{\sigma S}{2\sqrt{T}}\varphi(a_1)
+ rKe^{-rT}\Phi(a_1-\sigma\sqrt{T}).
\end{align}
\end{definition}

\begin{itemize}
\item (コール)オプションでは $\Theta > 0$ になる(満期が長いほど価値が高い
から)。
\item 通常，$\Theta$ は満期付近では小さく，満期の少し手前で上昇し，その後は
満期から離れるほど低下する。ただし，アットザマネー近くの場合，満期直前まで
$\Theta$ は大きい値をとる($K=100$， $r=0.1$， $\sigma=0.2$ の例で，$S=98, 100,
102$ の3ケースの比較が講義で示された)。
\item 満期 $T$ の変化は確定的なので，あらかじめ予測可能・対処可能である。
\end{itemize}

\subsection{ロー($\rho$)，ベガ($\mathcal{V}$)}

\begin{itemize}
\item \textbf{ロー} (Rho)：金利 $r$ が微小量変化したときの，オプション価格の
変化率。
\begin{equation}
\rho = \frac{\partial V_0^C}{\partial r}.
\end{equation}
\item \textbf{ベガ} (Vega)：ボラティリティ $\sigma$ が微小量変化したときの，
オプション価格の変化率。
\begin{equation}
\nu = \frac{\partial V_0^C}{\partial\sigma}.
\end{equation}
\end{itemize}

vega risk は，市場では非常に注目されている(市場では，ボラティリティ $\sigma$ は
一定でないと認識されているからである)。vega の変化，特に vanna
($\partial\nu/\partial S$)と volga($\partial\nu/\partial\sigma$)も注目
されている。

\section{通貨・金利デリバティブへの応用}

\subsection{通貨デリバティブへの応用}

為替レートも確率変動するとみなして，Black-Scholesモデルを適用するという
考え方は，とてもわかりやすい。素朴には
\[
\frac{\dd X(t)}{X(t)} = r(t)\,\dd t + \sigma\,\dd W(t)\ ?
\]
としたくなるが，それぞれの国に無リスク金利があること，裁定機会の排除を考えると
(この講義の水準を越えるので詳細は略)，正しくは
\begin{equation}
\frac{\dd X(t)}{X(t)} = \bigl(r_d(t) - r_f(t)\bigr)\dd t + \sigma\,\dd W(t)
\end{equation}
となる。ここで，添字 $d$ は基準通貨，$f$ は別の通貨，$X(t)$ は基準通貨への
換算レート，$B_i(t)$ は $i$ 通貨の無リスク資産の時刻 $t$ における価値である。
こういったモデルを基準にして，考えていくことができる。

\subsection{金利デリバティブへの応用}

金利も確率変動するとみなして，Black-Scholesモデルを適用するという考え方は，
とてもわかりやすい。では，それは可能か(金利 $=$ 無リスク金利として考える)。

\begin{itemize}
\item 金利自体が市場取引される資産ではない。金利は利子・利息計算や割引計算で
使用される\emph{指標}に過ぎない。
\item 金利自体が利回り・収益率である。そのため，金利自身の(瞬間的な)期待
収益率が(無リスク)金利自身($\mu = r(t)$)とするのはおかしい。
$\to$ この表現はやや不適切。
\end{itemize}
考え方は少し複雑であり，この講義の水準を越えるので略す。

% ---------------- 演習問題 ----------------
\section{演習問題}

以下の演習問題について，解答例のPDFは配布資料に含まれていなかったため，
解答・解説は編集者が作成したものである。

\begin{exercise}[問題12-1]
Black-Scholesの公式でプット・コール・パリティを確認しなさい。
\end{exercise}

\begin{answerbox}
\textbf{考え方。}BSのコール価格式\eqref{eq:bs-call}とプット価格式
\eqref{eq:bs-put}の差を計算し，$\Phi(x) + \Phi(-x) = 1$ を使う。

\medskip
\textbf{解答。}$a_2 = a_1 - \sigma\sqrt{T}$ と書くと，
\begin{align*}
V_0^C - V_0^P
&= \bigl[S\Phi(a_1) - Ke^{-rT}\Phi(a_2)\bigr]
- \bigl[Ke^{-rT}\Phi(-a_2) - S\Phi(-a_1)\bigr]\\
&= S\bigl[\Phi(a_1) + \Phi(-a_1)\bigr]
- Ke^{-rT}\bigl[\Phi(a_2) + \Phi(-a_2)\bigr]\\
&= S - Ke^{-rT}.
\end{align*}
金利一定のもとで $v(0,T) = e^{-rT}$ なので，これは
\[
V_0^P = V_0^C - S + K\,v(0,T),
\]
すなわちプット・コール・パリティそのものである。$\blacksquare$
\end{answerbox}

\begin{exercise}[問題12-2]
Black-Scholesモデルにおけるヨーロピアン・プット・オプションの価格公式
$V_0^P = V_0^P(S,K,T,\sigma,r)$ の振る舞いを示しなさい。
\end{exercise}

\begin{answerbox}
\textbf{考え方。}プット・コール・パリティ $V_0^P = V_0^C - S + Ke^{-rT}$ を
使えば，本文で示したコールの振る舞いから，プットの振る舞いが直ちに得られる。

\medskip
\textbf{解答。}
\begin{itemize}
\item 原資産価格 $S$：$S\to 0$ で $V_0^C\to 0$ より
$V_0^P \to Ke^{-rT}$(確実に行使され，行使価格の現在価値になる)。
$S\to\infty$ では，$V_0^C - S \to 0$ となるので $V_0^P \to 0$。
\item 行使価格 $K$：$K\to 0$ で $V_0^P \to S - S = 0$。$K\to\infty$ で
$V_0^P \approx Ke^{-rT} - S \to \infty$。
\item 満期 $T$：$T\to 0$ で
$V_0^P \to \max(S-K,0) - S + K = \max(K-S,\,0)$(本源的価値に収束)。
$T\to\infty$ では $V_0^C\to S$，$Ke^{-rT}\to 0$($r>0$)より $V_0^P \to 0$。
\item ボラティリティ $\sigma$：$\sigma\to 0$ で
$V_0^P \to \max(S-Ke^{-rT},0) - S + Ke^{-rT} = \max(Ke^{-rT}-S,\,0)$。
$\sigma\to\infty$ で $V_0^C\to S$ より $V_0^P \to Ke^{-rT}$
(プットの価格上限)。
\item 無リスク金利 $r$：$r\to 0$ で $V_0^P \to V_0^P(S,K,T,\sigma,0)$。
$r\to\infty$ で $V_0^C\to S$，$Ke^{-rT}\to 0$ より $V_0^P\to 0$。
\end{itemize}

\medskip
\textbf{ポイント。}コールの価値は $S$ を上限として無限に増え得るのに対し，
プットの価値は $Ke^{-rT}$ が上限である(受け取れる金額が最大でも $K$ だから)。
また，満期を延ばすとコールは単調に価値が上がるが，プットは
「受取 $K$ の割引が進む」効果があるため，必ずしも単調ではない(これが
ヨーロピアン・プットの時間価値が負になり得ることと対応している)。
\end{answerbox}

\begin{exercise}[問題12-3]
ブラックショールズの公式のうち，ヨーロピアン・コール・オプションの価格式から，
グリークスを求めなさい。
(1) デルタ delta，(2) ガンマ gamma，(3) セータ theta，(4) ロー rho，
(5) ヴェガ vega(or ラムダ lambda)。
\end{exercise}

\begin{answerbox}
\textbf{考え方。}コール価格式
$V_0^C = S\Phi(a_1) - Ke^{-rT}\Phi(a_2)$，\ $a_2 = a_1 - \sigma\sqrt{T}$，
\[
a_1 = \frac{\log(S/K) + (r+\sigma^2/2)T}{\sigma\sqrt{T}}
\]
を各変数で偏微分する。計算を簡潔にする鍵は，標準正規密度
$\varphi(x) = \Phi'(x) = \frac{1}{\sqrt{2\pi}}e^{-x^2/2}$ に関する恒等式
\begin{equation}
S\,\varphi(a_1) = Ke^{-rT}\varphi(a_2)
\tag{$\ast$}
\end{equation}
である。($\ast$)は
$\varphi(a_2) = \varphi(a_1)e^{a_1\sigma\sqrt{T} - \sigma^2T/2}$ と
$e^{a_1\sigma\sqrt{T}} = (S/K)e^{(r+\sigma^2/2)T}$ を代入すれば確かめられる。

\medskip
\textbf{(1) デルタ。}$a_1, a_2$ の $S$ 依存性に注意して
\[
\Delta = \frac{\partial V_0^C}{\partial S}
= \Phi(a_1) + S\varphi(a_1)\frac{\partial a_1}{\partial S}
- Ke^{-rT}\varphi(a_2)\frac{\partial a_2}{\partial S}.
\]
$\partial a_1/\partial S = \partial a_2/\partial S = 1/(S\sigma\sqrt{T})$ と
($\ast$)より後ろの2項は相殺し，
\[
\boxed{\Delta = \Phi(a_1)}.
\]

\textbf{(2) ガンマ。}
\[
\Gamma = \frac{\partial\Delta}{\partial S}
= \varphi(a_1)\frac{\partial a_1}{\partial S}
= \boxed{\frac{\varphi(a_1)}{S\sigma\sqrt{T}}} > 0.
\]

\textbf{(3) セータ。}$T$ で偏微分する。
\[
\Theta = \frac{\partial V_0^C}{\partial T}
= S\varphi(a_1)\frac{\partial a_1}{\partial T}
+ rKe^{-rT}\Phi(a_2) - Ke^{-rT}\varphi(a_2)\frac{\partial a_2}{\partial T}.
\]
($\ast$)を使うと第1項と第3項は
$S\varphi(a_1)\bigl(\partial a_1/\partial T - \partial a_2/\partial T\bigr)$ に
まとまり，$a_1 - a_2 = \sigma\sqrt{T}$ より
$\partial(a_1-a_2)/\partial T = \sigma/(2\sqrt{T})$ なので，
\[
\boxed{\Theta = \frac{\sigma S}{2\sqrt{T}}\,\varphi(a_1)
+ rKe^{-rT}\,\Phi(a_2)} > 0.
\]
(本文のセータの式と一致する。)

\textbf{(4) ロー。}
\[
\rho = \frac{\partial V_0^C}{\partial r}
= S\varphi(a_1)\frac{\partial a_1}{\partial r}
+ KTe^{-rT}\Phi(a_2) - Ke^{-rT}\varphi(a_2)\frac{\partial a_2}{\partial r}.
\]
$\partial a_1/\partial r = \partial a_2/\partial r = \sqrt{T}/\sigma$ と($\ast$)
より第1項と第3項は相殺し，
\[
\boxed{\rho = KTe^{-rT}\,\Phi(a_2)} > 0.
\]

\textbf{(5) ベガ。}
\[
\nu = \frac{\partial V_0^C}{\partial\sigma}
= S\varphi(a_1)\frac{\partial a_1}{\partial\sigma}
- Ke^{-rT}\varphi(a_2)\frac{\partial a_2}{\partial\sigma}
= S\varphi(a_1)\Bigl(\frac{\partial a_1}{\partial\sigma}
- \frac{\partial a_2}{\partial\sigma}\Bigr)
\]
(2つ目の等号で($\ast$)を使用)。$a_1 - a_2 = \sigma\sqrt{T}$ より
$\partial(a_1-a_2)/\partial\sigma = \sqrt{T}$ なので，
\[
\boxed{\nu = S\sqrt{T}\,\varphi(a_1)} > 0.
\]

\medskip
\textbf{ポイント。}どの計算でも，恒等式($\ast$)によって「$a_1,a_2$ の中身の
微分から出てくる項」がすべて相殺される。結果として，デルタは行使確率に似た
$\Phi(a_1)$，ガンマとベガはATM付近で最大になる $\varphi(a_1)$ 型，ローと
セータの主要項は行使確率 $\Phi(a_2)$ 型，というきれいな構造になる。符号は
コールではすべて正である($\Theta>0$ は「満期までの期間 $T$ が長いほど価値が
高い」ことを表す)。
\end{answerbox}

\section*{講義のまとめにかえて}

本講義では，金融市場の基本(第1〜3章)，リスクと効用の考え方(第4章)，
現代ポートフォリオ理論(第5〜7章)，デリバティブとその価格付け(第8〜12章)を
学んだ。価格付けの中心にあるのは「無裁定」と「複製」，そしてそれらから導かれる
「リスク中立化法」である。講義の構成に挙げられていた「9. 金融リスク管理の基礎」
(リスク尺度，市場リスクと信用リスク，金融危機と金融規制)は，第4章の
VaR・ESへの言及，第8章のCDS，第11章付録の信用リスク評価で，その入口に触れて
いる。
